\section{A Ramsey Model with Natural Resource Extraction and Recycling of Materials: Tastes and Technologies}
\label{sec:model-setup}

We consider an economy inhabited by a representative family dynasty with an infinite horizon. In each period the family derives utility $u(C)$ from consumption of final goods $C$ and incurs disutility $v(P)$ from the stock of pollution $P$ accumulated in the environment. The marginal utility of consumption is positive but declining, while the marginal disutility from pollution is positive and increasing, so at time zero the present value of the family's lifetime utility is

\begin{align}
U_0 &= \int_0^{\infty} \left[u(C) - v(P)\right] e^{-\rho t} dt, \quad  u'>0, \quad u''<0, \quad v'>0 \quad v''<0,
\label{eq:ramsey:utility}
\end{align}
where $\rho > 0$ is the constant rate of time preference, and the variables $C$ and $P$ are understood to be functions of time $t$.

The output of final goods $Y$ used for consumption and investment varies positively with the inputs of capital $K^Y$ and raw materials $R$, with decreasing marginal products of both inputs, while pollution harms productivity. Hence we have
\begin{align}
Y = F(K^Y, R, P),
\label{eq:ramsey:output}
\end{align}
where 
\begin{align*}
F_K &\equiv \frac{\partial F}{\partial K^Y} > 0, &
F_{KK} &\equiv \frac{\partial^2 F}{(\partial K^Y)^2} < 0, & F_R &\equiv \frac{\partial F}{\partial R} > 0,
&
F_{RR} &\equiv \frac{\partial^2 F}{\partial R^2} < 0,
& F_P &\equiv \frac{\partial F}{\partial P} < 0.
\end{align*}

The total input of raw materials in final goods production is the sum of virgin materials $N$ newly extracted from nature and recycled materials $R^R$:
\begin{align}
R &= N + R^R.
\label{eq:ramsey:materials}
\end{align}
In the recycling process an amount $\varpi W$ of waste products from production and consumption processes is transformed into usable raw materials via the technology
\begin{align}
R^R &= a\left(\frac{K^R}{\varpi W}\right) \varpi W,
\label{eq:ramsey:recycling} 
\end{align}
where 
\begin{align*}
a(0)= 0, && a'> 0, && a'' < 0, && \varepsilon \equiv a' \frac{K^R/(\varpi W)}{a} < 1, && \lim_{K^R/(\varpi W) \to \infty} a = 1,
\end{align*}
where $K^R$ is capital invested in recycling, $W$ is the total quantity of waste created in production and consumption, and $\varpi$ is the fraction of waste which is treated to make it amenable to recycling. The function $a\left(\frac{K^R}{\varpi W}\right)$ is the recycling rate with an elasticity $\varepsilon < 1$ with respect to the capital intensity $K^R/(\varpi W)$, reflecting diminishing returns in the recycling process. The assumption $a(0)=0$ implies that no recycling is possible if no capital is invested in recycling equipment. According to the last assumption in \eqref{eq:ramsey:recycling}, a complete recycling of all waste products $(a=1)$ would require an infinitely high capital intensity of the recycling process and is thus infeasible in practice, in line with the argument made by \textcite{georgescu1971} that some matter is inevitably degraded in the economic process and becomes unavailable for human purposes.\footnote{Georgescu-Roegen argued that the impossibility of 100 percent recycling follows from the second law of thermodynamics (the entropy law) which -- according to him -- applies to matter as well as energy. Many other scholars have pointed out that the laws of thermodynamics do not rule out the possibility of complete recycling of matter if only enough energy is available, but there is almost universal agreement that the required amounts of energy would be so large that the resulting pressure on the environment would be unsustainable. See \textcite{Couix2019} for an overview of this debate.} The most straightforward interpretation of the recycling technology in \eqref{eq:ramsey:recycling} is that it applies to minerals, but it may also encompass important forms of abatement of greenhouse gases. For example, the \CO{2} generated by the burning of fossil fuel may be captured and used as an input in certain forms of production via technologies for Carbon Capture and Utilization (CCU), and there are also technologies for capture and utilization of the \CO{2} and methane contained in the biogenic waste products generated by agriculture.

The total stock of man-made capital $K$ used to produce final goods and to recycle materials is
\begin{align}
K &= K^Y + K^R.
\label{eq:ramsey:capital-allocation}
\end{align}
We may think of the capital stocks as composites of physical and human capital where optimizing behaviour ensures that investments in the two forms of capital yield the same marginal return.

Virgin raw materials are extracted from a stock $S$ of known reserves of an exhaustible natural resource. The costs per period of extracting the flow $N$, measured in units of the final good and denoted by $C^N$, are given by the following cost function displaying positive and non-decreasing marginal extraction costs:\footnote{Working directly with a reduced-form cost function rather than deriving it from explicit assumptions about the underlying technology simplifies the exposition without altering our qualitative results.}
\begin{align}\label{eq:ramsey:extraction-cost}
C^N = C^N(N,S), && C_N^N \equiv \frac{\partial C^N}{\partial N} > 0, && C_{NN}^N \equiv \frac{\partial^2 C^N}{\partial N^2} \ge 0, &&
C_S^N \equiv \frac{\partial C^N}{\partial S} < 0.
\end{align}
According to \eqref{eq:ramsey:extraction-cost}, a larger stock of reserves facilitates extraction, thereby reducing extraction costs.

The stock of known reserves of the natural resource may be augmented through exploration activity which generates the total cost
\begin{align}\label{eq:ramsey:discovery-cost}
C^D &= C^D(D,X), &&
C_D^D \equiv \frac{\partial C^D}{\partial D} > 0,
&&
C_{DD}^D \equiv \frac{\partial^2 C^D}{\partial D^2} \ge 0,
&&
C_X^D \equiv \frac{\partial C^D}{\partial X} > 0,
\end{align}
where $D$ is the volume of new reserves discovered during the period considered, and $X$ is the total accumulated volume of reserves discovered in the past. According to \eqref{eq:ramsey:discovery-cost}, the marginal cost of discovery is positive and non-decreasing. The assumption $C_X^D > 0$ reflects that new reserves become harder to find the larger the total reserves already discovered.

The total amount of waste generated per period is proportional to the various economic activities:
\begin{align}\label{eq:ramsey:waste-generation}
W = \Omega\left(\omega^Y Y + \omega^N N + \omega^D D + \omega^C C\right), \qquad \Omega > 0.
\end{align}
The variable $\Omega$ reflects the average waste intensity of the various economic activities, while the positive $\omega$-coefficients in \eqref{eq:ramsey:waste-generation} capture possible differences in waste intensity across activities.

We assume that all waste has to be collected and that waste delivered for recycling must undergo some treatment to become recyclable. The collection cost per unit of waste collected is $c^c$, and the treatment cost per unit of total waste is $c^T$, so the total cost of waste collection and treatment, $C^W$, is:\footnote{A straightforward generalization of \eqref{eq:ramsey:waste-cost} would be to assume that there is a separate unit cost of collecting and treating each of the different types of waste in \eqref{eq:ramsey:waste-generation}. Similarly, we could assume that there is a separate recycling technology of the general form \eqref{eq:ramsey:recycling} for each category of waste. Such extensions would not add anything of substance to our analysis.}
\begin{align}\label{eq:ramsey:waste-cost}
C^W = \left[c^c + c^T(\varpi)\right] W, \qquad c^T(0) = \left.\frac{dc^T}{d\varpi}\right|_{\varpi=0} = 0, \qquad \frac{dc^T}{d\varpi} > 0, \quad \frac{d^2c^T}{d\varpi^2} > 0 \quad \text{for } \varpi > 0.
\end{align}
According to \eqref{eq:ramsey:waste-cost}, the marginal unit cost of treating the waste increases with the fraction $\varpi$ of waste treated.

Denoting total net investment in capital by $I$ and the exponential rate of depreciation of the existing capital stock by $\delta$, the economy's aggregate resource constraint expressed in units of the final good may now be stated as
\begin{align}
Y &= C + I + \delta K + C^N + C^D + C^W,
\label{eq:ramsey:resource-constraint}
\end{align}
and with a dot above a variable indicating its derivative with respect to time, the change in the total capital stock over time is
\begin{align}
\dot{K} &= I.
\label{eq:ramsey:capital-motion}
\end{align}

In each period the stock of known reserves of the natural resource is augmented by that period's new discoveries and is reduced by the resources extracted from existing reserves,
\begin{align}
\dot{S} &= D - N,
\label{eq:ramsey:reserves-motion}
\end{align}
and during each period the accumulated discoveries increase by that period's new discoveries:
\begin{align}
\dot{X} &= D.
\label{eq:ramsey:discoveries-motion}
\end{align}

The environment has some capacity to absorb and dissipate waste products which varies negatively with the stock of pollution already accumulated, so ceteris paribus the stock of pollution falls by the amount $\theta(P)P$ in each period via natural processes, where $\theta'(P) < 0$. At the same time, by normalization the stock of pollution increases by the amount of untreated waste $(1-\varpi)W$ directly deposited in or emitted to the environment. The amount of treated waste that cannot be recycled, $\varpi(1-a)W$, is assumed to augment the stock of pollution by the amount $d^W \varpi(1-a)W$, so the pollution stock evolves as
\begin{align}
\dot{P} &= \left[1 - \varpi + d^W\varpi\left(1 - a\left(\frac{K^R}{\varpi W}\right)\right)\right] W - \theta(P)P, \qquad 0 < d^W < 1, \qquad \theta'(P)<0.
\label{eq:ramsey:pollution-motion}
\end{align}
The assumption $d^W < 1$ reflects that treated non-recyclable waste is less harmful to the environment than untreated waste.

\subsection{The Materials Balance Principle}
\label{subsec:model-setup:materials-balance}

The materials balance principle implies that all materials extracted from nature ultimately end up as a similar quantity of waste of some kind, including various forms of emissions. However, the process of waste generation may be delayed in so far as the extracted materials are temporarily stored in accumulated stocks of durable goods, so it is only the current depreciation of the capital stock (due to wear and tear) that contributes to the current period's flow of waste products. As indicated in \eqref{eq:ramsey:waste-generation}, the total amount of waste generated by the mining sector during the current period is $\Omega\omega^N N + \Omega\omega^D D$. This waste represents additional materials that have to be extracted from the environment in order to obtain an amount $N$ of raw material that can be used as an input in final goods production. For example, to obtain a quantity $N$ of useful iron, it is necessary to extract a quantity $N(1+\Omega\omega^N)$ of iron ore, and in order to locate the iron ore, it is necessary to ``remove'' some material from the crust of the earth, which generates an amount of waste $\Omega\omega^D D$ in the process of exploration. In each period, the materials balance principle therefore implies that
\begin{align}
N\left(1+\Omega\omega^N\right) + \Omega\omega^D D &= W + \phi^K \dot{K},
\label{eq:model-setup:materials-balance}
\end{align}
where $\phi^K$ is the material intensity of capital goods, and where the real capital stock $K$ is measured in value terms in fixed prices (or chained values) as in the national accounts, whereas $N$ and $D$ and the associated waste flows are measured in tonnes. The left side of \eqref{eq:model-setup:materials-balance} measures the total mass of material extracted from nature in the period considered. According to the materials balance principle, this material must either be stored in durable goods via the current period's net accumulation of capital, $\dot K$, or it must end up in some kind of waste, $W$. Inserting \eqref{eq:ramsey:waste-generation} and \eqref{eq:ramsey:capital-motion} in \eqref{eq:model-setup:materials-balance}, we get
\begin{align}
N &= \phi^K I + \Omega\left(\omega^Y Y + \omega^C C\right).
\label{eq:model-setup:materials-balance-reduced}
\end{align}
If the left side of \eqref{eq:model-setup:materials-balance} is reinterpreted to represent the total amount of energy used in production and consumption, and the right side is redefined to represent the various forms taken by the energy after its use in economic processes, with $W$ including waste heat and emissions from energy use, and $\phi^K I$ capturing any energy temporarily stored in some form of capital, \eqref{eq:model-setup:materials-balance} may also be seen as a reflection of the first law of thermodynamics, according to which energy can neither be created nor destroyed.

Equation~\eqref{eq:model-setup:materials-balance-reduced} is the version of the materials balance principle that will be incorporated as a constraint in our model of the circular economy. We will treat the average waste intensity $\Omega$ and the relative waste intensities $\omega^N$, $\omega^D$, $\omega^Y$, and $\omega^C$ as exogenous variables determined by production and consumption technologies.

This completes the description of tastes and technologies in the economy. To simplify the exposition, we have not stated explicitly that the production function $F(\cdot)$, the cost functions $C^N(\cdot)$, $C^D(\cdot)$, and $C^W(\cdot)$, as well as the waste coefficients in \eqref{eq:ramsey:waste-generation} may depend on time as a result of technical and structural change, but this will be clarified in Section~\ref{sec:quantitative-analysis}, which presents the numerical version of the model. Before then we will characterize the first-best optimal allocation of resources and analyse how that allocation could be implemented in a market economy.
